% Created 2024-12-10 mar 16:19
% Intended LaTeX compiler: pdflatex
\documentclass[presentation]{beamer}
\usepackage[utf8]{inputenc}
\usepackage[T1]{fontenc}
\usepackage{graphicx}
\usepackage{longtable}
\usepackage{wrapfig}
\usepackage{rotating}
\usepackage[normalem]{ulem}
\usepackage{amsmath}
\usepackage{amssymb}
\usepackage{capt-of}
\usepackage{hyperref}
\usepackage{minted}
\setbeamertemplate{headline}{}
\usefonttheme{serif}
\usetheme{metropolis}
\usetheme{default}
\author{Tu Nombre}
\date{\today{}}
\title{Bases de Datos NoSQL: MongoDB y Cassandra}
\hypersetup{
 pdfauthor={Tu Nombre},
 pdftitle={Bases de Datos NoSQL: MongoDB y Cassandra},
 pdfkeywords={},
 pdfsubject={},
 pdfcreator={Emacs 27.1 (Org mode 9.6.18)}, 
 pdflang={English}}
\begin{document}

\maketitle

\begin{frame}[label={sec:org621383a}]{Introducción}
Las bases de datos NoSQL han ganado popularidad en Big Data:
\begin{itemize}
\item Manejan datos no estructurados y semi-estructurados.
\item Ofrecen escalabilidad horizontal.
\end{itemize}

En esta presentación, profundizaremos en:
\begin{itemize}
\item \alert{\alert{MongoDB}}.
\item \alert{\alert{Cassandra}}.
\end{itemize}

---
\end{frame}

\begin{frame}[label={sec:org2994fa3}]{MongoDB: Base de Datos Documental}
\begin{block}{Características Clave}
\begin{enumerate}
\item \alert{\alert{Modelo de Documentos}}:
\begin{itemize}
\item Estructuras en JSON o BSON.
\item Datos complejos y anidados sin esquema fijo.
\end{itemize}
\item \alert{\alert{Escalabilidad Horizontal}}:
\begin{itemize}
\item Soporta \alert{\alert{sharding}} para distribuir datos.
\end{itemize}
\item \alert{\alert{Consulta Flexible}}:
\begin{itemize}
\item Operaciones avanzadas como agregaciones y búsquedas geoespaciales.
\end{itemize}
\item \alert{\alert{Alta Disponibilidad}}:
\begin{itemize}
\item Réplicas con Replica Sets.
\end{itemize}
\item \alert{\alert{Compatibilidad Multiplataforma}}:
\begin{itemize}
\item Integración con aplicaciones modernas.
\end{itemize}
\end{enumerate}
\end{block}

\begin{block}{Ventajas}
\begin{itemize}
\item Flexibilidad con datos dinámicos.
\item Rendimiento alto en escrituras.
\item Estructura intuitiva para desarrolladores.
\end{itemize}
\end{block}

\begin{block}{Desventajas}
\begin{itemize}
\item Consistencia eventual en entornos distribuidos.
\item Mayor consumo de almacenamiento con BSON.
\end{itemize}
\end{block}

\begin{block}{Casos de Uso}
\begin{itemize}
\item Aplicaciones web y móviles dinámicas.
\item Gestión de catálogos de productos.
\item Búsquedas avanzadas geoespaciales.
\end{itemize}

---
\end{block}
\end{frame}

\begin{frame}[label={sec:org652cd06},fragile]{Ejemplo Práctico: MongoDB}
 \begin{block}{Gestión de un Catálogo de Productos}
Ejemplo de documento en JSON:
\begin{verbatim}
{
  "productId": "12345",
  "name": "Laptop",
  "brand": "TechBrand",
  "specifications": {
    "processor": "Intel i7",
    "RAM": "16GB",
    "storage": "512GB SSD"
  },
  "price": 1200,
  "categories": ["Electronics", "Computers"]
}
\end{verbatim}

---
\end{block}
\end{frame}

\begin{frame}[label={sec:org43493f4}]{Cassandra: Base de Datos de Columnas Anchas}
\begin{block}{Características Clave}
\begin{enumerate}
\item \alert{\alert{Modelo de Columnas Anchas}}:
\begin{itemize}
\item Organiza datos en filas y columnas con número variable de columnas.
\end{itemize}
\item \alert{\alert{Alta Escalabilidad Horizontal}}:
\begin{itemize}
\item Escala hasta miles de nodos sin impacto en rendimiento.
\end{itemize}
\item \alert{\alert{Consistencia Configurable}}:
\begin{itemize}
\item Modelo CAP: consistencia fuerte o eventual.
\end{itemize}
\item \alert{\alert{Replicación Inteligente}}:
\begin{itemize}
\item Garantiza tolerancia a fallos y alta disponibilidad.
\end{itemize}
\item \alert{\alert{Altamente Disponible}}:
\begin{itemize}
\item Sin puntos únicos de fallo.
\end{itemize}
\end{enumerate}
\end{block}

\begin{block}{Ventajas}
\begin{itemize}
\item Escalabilidad para grandes volúmenes de datos.
\item Tolerancia a fallos integrada.
\item Rendimiento óptimo en escrituras masivas.
\end{itemize}
\end{block}

\begin{block}{Desventajas}
\begin{itemize}
\item Complejidad en configuración.
\item Operaciones limitadas comparadas con SQL.
\end{itemize}
\end{block}

\begin{block}{Casos de Uso}
\begin{itemize}
\item Monitoreo de IoT.
\item Gestión de logs distribuidos.
\item Análisis de datos financieros.
\end{itemize}

---
\end{block}
\end{frame}

\begin{frame}[label={sec:org6219490},fragile]{Ejemplo Práctico: Cassandra}
 \begin{block}{Registro de Logs de Servidores}
Modelo de tabla en CQL:
\begin{minted}[]{sql}

CREATE TABLE server_logs (
  log_id UUID PRIMARY KEY,
  timestamp TIMESTAMP,
  server_id TEXT,
  log_message TEXT
);
\end{minted}

---
\end{block}
\end{frame}

\begin{frame}[label={sec:orgda32e5c}]{Comparación: MongoDB vs Cassandra}
\end{frame}
\begin{frame}[label={sec:orga2733ac}]{Comparación: MongoDB vs Cassandra}
\begin{center}
\begin{tabular}{lll}
\alert{\alert{Aspecto}} & \alert{\alert{MongoDB}} & \alert{\alert{Cassandra}}\\[0pt]
\hline
Modelo de Datos & Documentos (JSON/BSON) & Columnas anchas (Key-Value extendido)\\[0pt]
Escalabilidad & Alta, mediante sharding & Muy alta, ideal para datos distribuidos\\[0pt]
Consistencia & Eventual, con ACID parcial & Configurable: eventual o fuerte\\[0pt]
Rendimiento & Bueno para consultas complejas & Excelente para escrituras masivas\\[0pt]
Casos de Uso & Aplicaciones dinámicas, catálogos & IoT, logs distribuidos\\[0pt]
\end{tabular}
\end{center}

---
\end{frame}

\begin{frame}[label={sec:org2798ed1}]{Conclusión}
\begin{itemize}
\item \alert{\alert{MongoDB}}: Ideal para aplicaciones con datos dinámicos y análisis flexible.
\item \alert{\alert{Cassandra}}: Perfecta para entornos distribuidos con requisitos de escritura intensiva.
\item La elección depende del caso de uso y los requisitos específicos del proyecto.
\end{itemize}
\end{frame}
\end{document}
